\documentclass{article}
\usepackage[utf8]{inputenc}
\usepackage[russian]{babel}
\usepackage{amsmath}
\usepackage{amssymb}
\usepackage{mathtools}
\usepackage[a4paper, left=1.5cm, right=1.5cm, top=1.5cm, bottom=1.5cm]{geometry}\setlength{\parindent}{0pt}

\title{ Билет 10. Свойства интегралов, зависящих от параметра, в случае, когда пределы интегрирования также зависят от параметра.}
\author{}
\date{}

\begin{document}

\maketitle

Пусть задана функция $ f(x, y) $, определённая на прямоугольнике $ [a, b] \times [c, d] $. 

Пусть также заданы функции $ \alpha(y) $ и $ \beta(y) $, определённые при $ y \in [c, d] $, причём для всех $ y \in [c, d] $ выполняется:

\[
\alpha(y), \beta(y) \in [a, b]
\]

Определим функцию $ G(y) $ как интеграл с переменными пределами:

\[
G(y) = \int_{\alpha(y)}^{\beta(y)} f(x, y) \, dx
\]

*Рисунок.

\section*{Теорема о непрерывности по параметру}

Пусть $ f(x, y) $ — непрерывная функция двух переменных на $ [a, b] \times [c, d] $.

Тогда для любого $ y_0 \in [c, d] $:

\[
\lim_{y \to y_0} G(y) = G(y_0)
\]

Доказательство:

Рассмотрим разность:

\[
G(y) - G(y_0) = \int_{\alpha(y)}^{\beta(y)} f(x, y) \, dx - \int_{\alpha(y_0)}^{\beta(y_0)} f(x, y_0) \, dx
\]

Представим её как сумму трёх слагаемых:

\[
G(y) = \int_{\alpha(y_0)}^{\beta(y_0)} f(x, y) \, dx + \int_{\alpha(y)}^{\alpha(y_0)} f(x, y) \, dx + \int_{\beta(y_0)}^{\beta(y)} f(x, y) \, dx
\]

Первое слагаемое: 
\[
\int_{\alpha(y_0)}^{\beta(y_0)} f(x, y) \, dx \to \int_{\alpha(y_0)}^{\beta(y_0)} f(x, y_0) \, dx = G(y_0)
\quad \text{при } y \to y_0,
\]
поскольку $ f(x, y) $ непрерывна по $ y $ равномерно по $ x $ (как непрерывная на прямоугольнике).

Второе и третье слагаемые:

\[
\left| \int_{\alpha(y)}^{\alpha(y_0)} f(x, y) \, dx \right| \leq M \cdot |\alpha(y) - \alpha(y_0)| \to 0 \quad \text{при } y \to y_0,
\]
если $ \alpha(y) $ непрерывна в точке $ y_0 $, и аналогично для $ \beta(y) $, где $ |f(x,y)| \leq M $ на прямоугольнике.\\
(Так как функция непрерывна на замкнутом ограниченном множестве, то по теореме Вейерштрасса она ограничена (и достигает своего максимума и минимума)).


Таким образом,
\[
\lim_{y\to y_0} \bigl(G(y)-G(y_0)\bigr)=0,
\]
то есть
\[
\lim_{y\to y_0} G(y)=G(y_0).
\]
Следовательно, функция $G(y)$ непрерывна в точке $y_0$.

\section*{Теорема о дифференцируемости по параметру}

Пусть:
\begin{itemize}
    \item $ f(x, y) $ определена и непрерывна на $ [a, b] \times [c, d] $,
    \item $ f_y'(x, y) $ существует и непрерывна на $ [a, b] \times [c, d] $,
    \item $ \alpha(y) $, $ \beta(y) $ непрерывно дифференцируемы на $ [c, d] $,
    \item $ \alpha(y), \beta(y) \in [a, b] $ для всех $ y \in [c, d] $.
\end{itemize}

Тогда $ G(y) $ дифференцируема на $ [c, d] $, и её производная вычисляется по формуле Лейбница:

\[
\exists G'(y_0) = \int_{\alpha(y_0)}^{\beta(y_0)} f_y'(x, y) \, dx - \alpha'(y_0) f(\alpha(y_0), y_0) + \beta'(y_0) f(\beta(y_0), y_0)
\]

\textbf{Доказательство:}

Произвольным образом выбираем $\forall y_0 \in [c, d]$. Рассматриваем $G(y)$ и снова записываем ее в форме

\[
\int_{\alpha(y_0)}^{\beta(y_0)} f(x, y) \, dx + \int_{\beta(y_0)}^{\beta(y)} f(x, y) \, dx + \int_{\alpha(y)}^{\alpha(y_0)} f(x, y) \, dx
\]

Ищем производную от этой функции. Первое слагаемое - пределы не зависят от $y$ и выполнены все условия о дифференцируемости под знаком интеграла. Значит у первого слагаемого производная есть, и она равна $\int_{\alpha(y_0)}^{\beta(y_0)} f_y'(x, y) \, dx$. Для второго слагаемого:
Пусть
\[
\Psi(y) = \int_{\beta(y_0)}^{\beta(y)} f(x, y) \, dx,
\]
где \(y_0\) — фиксированная точка, в которой ищем производную.

\textbf{Шаг 1. Определение производной:}
\[
\Psi'(y_0) = \lim_{y \to y_0} \frac{\Psi(y) - \Psi(y_0)}{y - y_0}.
\]
Так как \(\Psi(y_0) = \int_{\beta(y_0)}^{\beta(y_0)} f(x, y_0) dx = 0\), имеем:
\[
\Psi'(y_0) = \lim_{y \to y_0} \frac{1}{y - y_0} \int_{\beta(y_0)}^{\beta(y)} f(x, y) \, dx.
\]

\textbf{Шаг 2. Теорема о среднем значении для интеграла:} \\
\textit{Если функция \(f(x)\) непрерывна на отрезке \([p, q]\), то существует точка \(\xi \in [p, q]\), такая что}
\[
\int_p^q f(x) \, dx = f(\xi) \cdot (q - p).
\]

Применяем эту теорему к интегралу \(\int_{\beta(y_0)}^{\beta(y)} f(x, y) \, dx\) для каждого фиксированного \(y\). Получаем:
\[
\int_{\beta(y_0)}^{\beta(y)} f(x, y) \, dx = f(\xi_y, y) \cdot (\beta(y) - \beta(y_0)),
\]
где \(\xi_y \in [\beta(y_0), \beta(y)]\).

\textbf{Шаг 3. Подставляем в предел определения производной:}
\[
\Psi'(y_0) = \lim_{y \to y_0} \frac{f(\xi_y, y) \cdot (\beta(y) - \beta(y_0))}{y - y_0} = \lim_{y \to y_0} f(\xi_y, y) \cdot \frac{\beta(y) - \beta(y_0)}{y - y_0}.
\]

\textbf{Шаг 4. Переход к пределу:} \\
Так как \(f(x, y)\) непрерывна и \(\xi_y \to \beta(y_0)\) при \(y \to y_0\), получаем:
\[
\lim_{y \to y_0} f(\xi_y, y) = f(\beta(y_0), y_0).
\]

Кроме того,
\[
\lim_{y \to y_0} \frac{\beta(y) - \beta(y_0)}{y - y_0} = \beta'(y_0).
\]

Следовательно,
\[
\Psi'(y_0) = f(\beta(y_0), y_0) \cdot \beta'(y_0).
\]

А в пределе мы получим $(\beta'(y_0)) f(\beta(y_0), y_0)$ (по определению производной первый множитель и по непрерывности $f$ второй множитель). Аналогичным образом получим производную от третьего слагаемого, и она будет равна $-\alpha'(y_0) f(\alpha(y_0), y_0)$ (минус вылезет из-за того, что нам нужно будет перевернуть пределы интегрирования, чтобы $\alpha(y)$ было наверху). Собирая получим:

\[
G'(y_0) = \int_{\alpha(y_0)}^{\beta(y_0)} f_y'(x, y) \, dx + \beta'(y_0) f(\beta(y_0), y_0) - \alpha'(y_0) f(\alpha(y_0), y_0)
\]

Теорема доказана.
\end{document}